% !TeX root = ../main.tex
% !TeX root = ../main.tex


\chapter{Fundamentals}
\label{ch:fundamentals}

This chapter provides the concepts shared by the three studies. It first distinguishes
the POI prediction tasks (Section~\ref{sec:fund:tasks}). It then traces mobility
representations from one-hot identifiers to Check2HGI
(Section~\ref{sec:fund:repr}), explains the main forms and risks of multitask learning
(Section~\ref{sec:fund:mtl}), and defines the data and evaluation protocol
(Section~\ref{sec:fund:eval}). Section~\ref{sec:fund:relevance} connects these
elements to the questions addressed in the article chapters.


\section{Point-of-interest prediction tasks}
\label{sec:fund:tasks}

This section establishes what each model predicts before reviewing how those
predictions are made. Location-based social networks (LBSNs) record mobility as
check-ins, each linking a user, a POI, and a timestamp. Their geographic and temporal
detail supports data-driven studies of cities~\cite{silva2019urbancomputing}. The
records also reveal recurring behavior: people often return to a small set of places
and make longer trips less frequently~\cite{cho2011gowalla}. An entropy analysis
estimated the potential predictability of an individual's next location at about 93
percent~\cite{song2010limits}. Because this estimate concerns next-location prediction
at a coarse spatial resolution, it shows that mobility contains learnable regularity
but does not provide a reference point for the category and region metrics defined in
Section~\ref{sec:fund:eval}.

\subsection{Task boundaries and notation}

Human-mobility models address distinct targets, from crowd-flow estimation to
next-location prediction~\cite{luca2021mobilitysurvey}. Because these targets are not
interchangeable, model comparisons require explicit task boundaries. This dissertation
studies three tasks: one static and two sequential.


\subsubsection{Check-ins and histories}

Let $\mathcal{U}$, $\mathcal{P}$, $\mathcal{C}$, and $\mathcal{R}$ denote the sets of
users, POIs, category classes, and region classes. Each POI $p\in\mathcal{P}$ carries a
category $c_p\in\mathcal{C}$ and lies in a region $r_p\in\mathcal{R}$.

\begin{definition}[Check-in]\label{def:fund:checkin}
The $i$th check-in of user $u$ is the tuple
\begin{equation}
    x_i=(u,p_i,t_i,c_i,r_i),
\end{equation}
where $p_i\in\mathcal{P}$ is the visited POI, $t_i$ is its timestamp,
$c_i=c_{p_i}$ is its category, and $r_i=r_{p_i}$ is its region.
\end{definition}

\begin{definition}[Check-in history]\label{def:fund:history}
The check-in history of length $\ell$ preceding check-in $x_i$ is the ordered sequence
\begin{equation}
    H_i=(x_{i-\ell},\ldots,x_{i-1}).
\end{equation}
For the two sequential tasks studied here (Definitions~\ref{def:fund:nextcat}
and~\ref{def:fund:nextreg}), the prediction target is $c_i$ for next-category
prediction and $r_i$ for next-region prediction. These labels belong to $x_i$, not
to the preceding check-ins in $H_i$: the categories and regions carried by the
visits in $H_i$ are observed input, and only the attributes of $x_i$ itself are
withheld.
\end{definition}

\subsubsection{Check-in and place embedding}

A task definition must also state what the model reads. Two representations recur
throughout this dissertation, and the three studies vary which one supplies the model
input.

\begin{definition}[Place embedding]\label{def:fund:placelevel}
A place embedding assigns one vector $\mathbf{e}_p$ to each POI $p$, so every check-in
at that POI enters the model with the same representation.
\end{definition}

\begin{definition}[Check-in-level representation]\label{def:fund:checkinlevel}
A check-in-level representation assigns one vector $\mathbf{e}_{x_i}$ to each check-in
$x_i$, so two visits to the same POI may enter the model with different
representations.
\end{definition}

\begin{definition}[Representation map]\label{def:fund:repmap}
A representation map assigns to each check-in a vector,
\begin{equation}
    \rho(x_i)\in\mathbb{R}^{d},
\end{equation}
and extends to histories elementwise,
\begin{equation}
    \rho(H_i)=\bigl(\rho(x_{i-\ell}),\ldots,\rho(x_{i-1})\bigr).
\end{equation}
\end{definition}

Every model of the sequential tasks in this dissertation reads $\rho(H_i)$ rather than
$H_i$ itself. Chapter~\ref{ch:cbic} defines the input at the place level,
$\rho(x_j)=\mathbf{e}_{p_j}$. Chapter~\ref{ch:courb} retains a vector tied to the visited
POI but enriches it with separate spatial, temporal, and categorical components.
Chapter~\ref{ch:mobiwac} moves the input to the check-in level. Its semantic stream reads
the per-visit vector $\rho(x_i)=\mathbf{e}_{x_i}$, while its spatial stream reads the
trained vector of the visit's region node. The task definitions remain fixed, while
each study uses its own form of $\rho$.

\subsubsection{The three experimental tasks}

Three definitions separate the targets. The first is static and reads a place; the
other two are sequential and read a history.

\begin{definition}[Category classification]\label{def:fund:catclf}
Category classification predicts the category of a POI from its representation:
\begin{equation}
    g_{\mathrm{cat}}(\mathbf{e}_p)\longrightarrow c_p.
\end{equation}
At evaluation time, POI $p$ is held out from the classifier-training fold, so
``unknown POI'' means unknown to the classifier, not necessarily absent from the graph
used to learn $\mathbf{e}_p$.
\end{definition}

\begin{definition}[Next-category prediction]\label{def:fund:nextcat}
Next-category prediction maps a check-in history to the category of the next visit:
\begin{equation}
    f_{\mathrm{cat}}(H_i)\longrightarrow c_i.
\end{equation}
\end{definition}

\begin{definition}[Next-region prediction]\label{def:fund:nextreg}
Next-region prediction maps a check-in history to the region of the next visit:
\begin{equation}
    f_{\mathrm{reg}}(H_i)\longrightarrow r_i.
\end{equation}
\end{definition}

The task pairing changes across the studies. Chapters~\ref{ch:cbic}
and~\ref{ch:courb} combine category classification
(Definition~\ref{def:fund:catclf}) with next-category prediction
(Definition~\ref{def:fund:nextcat}). Chapter~\ref{ch:mobiwac} combines next-category
prediction with next-region prediction (Definition~\ref{def:fund:nextreg}). All three
studies use seven categories: Community, Entertainment, Food, Nightlife, Outdoors,
Shopping, and Travel. For the region task, the target is an administrative unit at
neighborhood scale. These output spaces require different metrics and reference points,
and Section~\ref{sec:fund:eval} both names the unit each dataset supplies and explains
the measurement.

\subsection{Related sequential prediction settings}

One further target must be named, because it is the task that the models this section
reviews were built to solve.

\begin{definition}[Next-place prediction]\label{def:fund:nextplace}
Next-place prediction maps a check-in history to the identity of the next visited POI:
\begin{equation}
    f_{\mathrm{place}}(H_i)\longrightarrow p_i.
\end{equation}
It is named to delimit the scope of the dissertation, and no chapter reports a result
for $f_{\mathrm{place}}$.
\end{definition}

Next place is the dominant task in the field, and methods developed for it provide the
methodological basis for next-category and next-region prediction. Early approaches
relied on recurrent architectures. ST-RNN uses time-specific and distance-specific
transition matrices~\cite{liu2016strnn}, while DeepMove adds attention over long and
sparse visit histories~\cite{feng2018deepmove}. HST-LSTM incorporates
spatial-temporal context into a hierarchical recurrent cell~\cite{kong2018hstlstm}.
Flashback retrieves past hidden states according to the temporal and spatial gaps
between visits~\cite{yang2020flashback}.

Later approaches make attention central to the architecture. STAN connects
non-adjacent visits through spatio-temporal correlations~\cite{luo2021stan}. GeoSAN
combines a hierarchical geographic grid with self-attention~\cite{lian2020geosan}, and
GETNext introduces a global trajectory-flow map into a
transformer~\cite{yang2022getnext}. Every model named here predicts the exact next
place, so none of them is a direct baseline for the targets studied in this dissertation.

\subsection{Category and region as end targets}

Two strands of this literature support the task choice in this dissertation. The first
concerns how a visit is represented. CTLE learns context-aware and time-aware location
embeddings, allowing the same POI to have different vectors in different
visits~\cite{lin2021ctle}. This per-visit view motivates the representation discussed in
Section~\ref{sec:fund:repr}.

The second strand concerns the role of category and region predictions. Some methods
use them to support next-place prediction rather than as final outputs. HMT-GRN uses a
predicted region to constrain the search for a place~\cite{Lim2022}, and CatDM predicts
a category to reduce the set of candidate places~\cite{yu2020catdm}. Other studies make
one of these properties the final target. Activity-region prediction addresses the
region directly~\cite{zhu2022drrgnn}, while POI-RGNN predicts the next
category~\cite{capanema2023poirgnn}. This dissertation uses next category and next
region as final outputs of equal status in one model.


\section{Representations for mobility}
\label{sec:fund:repr}

This section traces the representations used in the dissertation and explains why the
final study moves from a place embedding to a check-in-level representation.

\subsection{From identifiers to graph embeddings}

The representation line begins with identifiers and progresses to vectors learned
from graph structure.

A one-hot identifier marks a place or category by position but encodes no relationship
between identifiers. Distributed representations instead learn dense vectors whose
geometry reflects relationships in the data. Skip-gram places symbols with similar
contexts near one another~\cite{mikolov2013word2vec}. DeepWalk transfers this idea to
graphs by treating random walks as sentences~\cite{perozzi2014deepwalk}, and node2vec
adjusts the walk to emphasize homophily or structural
roles~\cite{grover2016node2vec}. Graph neural networks learn these relationships
through neighborhood aggregation. Graph convolutional networks use a localized
spectral rule~\cite{kipf2017gcn}, graph attention networks learn the weight of each
neighbor~\cite{velivckovic2017graph}, and GraphSAGE learns aggregators that can embed
nodes not seen during training~\cite{hamilton2017graphsage}.

\subsection{Hierarchical graph infomax}

The graph-infomax methods supply the unsupervised representations used across the
three studies. The idea these methods share is simple enough to state before any of
the notation: the model learns useful vectors by being asked to tell a true pairing of
two parts of the data from a corrupted one, and it needs no labels to do so, because
the data itself says which pairing is the true one.

The representations used in this work are trained with no next-category and no
next-region target, each visit contributes with its own description: the category of
the place visited is an input feature of its node, together with the hour and the day
of the week.
Mutual-information estimators provide one way to learn such
representations~\cite{belghazi2018mine}. Deep InfoMax maximizes agreement between
local features and a global summary~\cite{hjelm2019dim}, and Deep Graph Infomax (DGI)
applies this objective to graph nodes and a graph-level
summary~\cite{velickovic2019deep}. DGI supplies the objective and the place embeddings
used in the first study. Hierarchical Graph Infomax (HGI) extends that objective
across POI, region, and city levels~\cite{huang2023hgi}.

HGI is the place-level representation that the final study replaces, so its reference
point matters here. The mechanism itself is set out in the original paper, and the
account below is limited to what the later chapters rely on. Its own stated goal is to
learn urban region representations from POIs in a fully unsupervised manner, and it
reaches that goal by extending the two scales of DGI to three: POI, region, and city.
The name states what is maximized. Training raises the mutual information between
representations at two adjacent levels of that hierarchy, and it does so without
evaluating that quantity in closed form. A bilinear discriminator combines two
embeddings through a learned weight matrix and passes the result through a logistic
function, and the loss rewards a high score for a true pair and a low score for a
false pair. No label of any downstream task enters that comparison, so the
representation is obtained without supervision.

One consequence of the design matters for the later chapters. A pretrained category
encoder supplies the initial POI features, one graph convolution layer over a Delaunay
graph of the POIs of the study area adds spatial context to each of them, multi-head
attention aggregates the POI embeddings of a region, and an area-weighted sum over
regions produces one city embedding. The training signal updates the POI encoder, the
aggregation function, and the region encoder together, so a POI embedding is optimized
to score high against the embedding of its own region and low against the embeddings
of other regions. Region membership also enters earlier, through the edge weights of
the POI graph, which decrease with the distance between two POIs and are reduced
further when the two POIs lie in different regions. The POI-level output is therefore
not a description of the place in isolation. It already reflects the region the place
belongs to.

HGI was developed and evaluated for urban region representation. Its reported
experiments estimate urban functional distributions, population density, and housing
price in Xiamen Island and Shenzhen, and the object under evaluation there is the
region embedding. The POI embedding is an internal stage in producing it. This
dissertation repurposes that POI-level output for sequential prediction, a use the
original evaluation does not cover. In this role, HGI supplies the place-level
baseline in the later studies and the direct basis of Check2HGI.

The move to other study areas left one setting of that baseline open. Huang et al.\
reduce an edge between POIs of different regions to 0.4 of its weight in the
experiments named above, and this work did not transfer that value untested. In
experiments on one of the state datasets of this dissertation, over the same five
folds used elsewhere, a value of 0.7 gave the best category F1 among the values tried,
and the later studies adopt it. What that comparison settles is one hyperparameter of
the place-level baseline. It is not evidence about how HGI and Check2HGI compare.

\subsection{From a fixed place vector to visit context}

The move from HGI to Check2HGI begins with information that a fixed POI vector cannot
express.

Definition~\ref{def:fund:placelevel} carries one limitation for this research: a
weekday morning and a Saturday night at the same place have identical inputs at the
representation level. CTLE shows an alternative by learning context-aware and
time-aware location embeddings whose vectors change with the
visit~\cite{lin2021ctle}.

Check2HGI reaches the check-in level of Definition~\ref{def:fund:checkinlevel} by
adding the check-in below the place in the hierarchy.

\subsubsection{Context encoders}

A per-visit representation needs temporal and spatial context in addition to the
identity of the visited POI. The second study supplies that context through separate
encoders while its input stays a function of the visited POI.
Time2Vec combines linear and periodic components for temporal
features~\cite{kazemi2019time2vec}. SIREN uses sinusoidal activation functions to
represent continuous signals~\cite{sitzmann2020implicit}. Space2Vec encodes location
and spatial context at several scales~\cite{mai2020multiscalerepresentationlearningspatial},
while Sphere2Vec represents coordinates on a sphere rather than on a flat
surface~\cite{mai2023sphere2vecgeneralpurposelocationrepresentation}. A related
approach combines spherical harmonics with sinusoidal
networks~\cite{russwurm2024geographiclocationencodingspherical}.
Chapter~\ref{ch:courb} keeps the MTLnet architecture but replaces its single
place embedding with spatial, temporal, and categorical encoders. The controlled
change isolates the effect of the input representation.

\subsubsection{The check-in level}


Check2HGI completes the move from contextualized inputs to one learned representation
per visit. It adds a fourth level below the place and learns one vector per visit from
the visits themselves, with no next-category and no next-region target in its
objective. The hierarchy has
three adjacent-level boundaries: check-in to place, place to region, and region to
city. Its three hierarchical boundaries contribute the fixed-weight sum
\begin{equation}\label{eq:fund:check2hgi}
    \mathcal{L} = 0.4\,\mathcal{L}_{c2p} + 0.3\,\mathcal{L}_{p2r} + 0.3\,\mathcal{L}_{r2c},
\end{equation}
where $\mathcal{L}_{c2p}$, $\mathcal{L}_{p2r}$, and $\mathcal{L}_{r2c}$ correspond
to these three boundaries. Each term uses a bilinear discriminator,
\begin{equation}\label{eq:fund:check2hgi-disc}
    \mathcal{D}(\mathbf{e}_1, \mathbf{e}_2) = \sigma(\mathbf{e}_1^{\top}\mathbf{W}\mathbf{e}_2),
\end{equation}
where $\mathbf{W}$ is learned and
$\sigma(z)=1/(1+e^{-z})$ is the logistic function. The discriminator is linear in
either embedding when the other is held fixed, and its output lies between 0 and 1.
For each boundary, the loss favors a high score for a true pair and a low score for a
false pair:
\begin{equation}\label{eq:fund:check2hgi-term}
    \mathcal{L}_{*} = -\log \mathcal{D}(\mathbf{e}^{+}, \mathbf{e}^{+})
    - \log\left(1 - \mathcal{D}(\mathbf{e}^{+}, \mathbf{e}^{-})\right),
\end{equation}
Here, $\mathbf{e}^{+}$ belongs to a true pair, such as a check-in and the place where
it occurred, whereas $\mathbf{e}^{-}$ is substituted from another example in the
batch. No next-step target appears in these equations. Two auxiliary terms complete the
objective that the final study trains: a reconstruction term at weight 0.3, which
recovers a masked place's own distribution over the seven categories from its
neighbors, and a term at weight 0.1, which keeps the trainable table of place vectors
near the values it started from.


Where a place embedding answers ``what is this place,'' a check-in-level representation
answers ``what is this visit,'' and
Chapter~\ref{ch:mobiwac} develops this representation and the joint model built on it.
Table~\ref{tab:fund:lineage} traces the full progression, from DGI through HGI to the
check-in level and the joint model that uses it.

\subsection{Model lineage}


Each study changes the input representation, and each also changes the architecture
that consumes it. Naming both sequences together lets a later chapter be read against
the step it altered.

The final joint model remains part of the MTLnet lineage. It retains the task-specific
input encoders, output heads, and division between shared and task-specific parameters.
MTLnet uses residual layers conditioned by Feature-wise Linear Modulation (FiLM) as its
shared middle. The joint model
replaces that component with cross-attention blocks, allowing each task stream to
read the other while retaining its own feed-forward parameters. A private spatial
path also gives the region output direct access to the region sequence before the
streams interact. The models therefore differ in their sharing topology and in the
private input available to the region output. Both architectures receive two inputs,
and what changes between them is the source: MTLnet derives both from one place
embedding, whereas the joint model reads two tables exported from the same check-in-level
representation. The two tables share an origin by construction, so they are not
independent views. This relationship allows the negative
result in Chapter~\ref{ch:cbic} and the positive result in
Chapter~\ref{ch:mobiwac} to be interpreted as two stages of one model lineage.

% !TeX root = ../../main.tex
\begin{table}[htbp]
  \centering
  \caption{Representation and model lineage threaded through this dissertation, from the
  place-level graph-infomax embeddings to the check-in-level representation and the joint
  model built on it. Status of the Check2HGI representation and the joint model is
  \emph{submitted, under review}.}
  \label{tab:fund:lineage}
  \begin{tabularx}{\textwidth}{@{}l>{\raggedright\arraybackslash}Xl@{}}
    \toprule
    Model / method & What it added & Reference \\
    \midrule
    DGI & Unsupervised place embeddings by maximizing mutual information between local patches and a global graph summary. & \cite{velickovic2019deep} \\
    \addlinespace
    HGI & Extends graph infomax over the POI, region, and city hierarchy to yield region-aware place embeddings. & \cite{huang2023hgi} \\
    \addlinespace
    MTLnet & First joint model: place embedding, FiLM conditioning, and hard parameter sharing. Null result for that configuration. & \cite{silva2025mtlnet} \\
    \addlinespace
    ST-MTLNet & Keeps MTLnet, replaces the place-embedding input with decomposed spatial, temporal, and categorical encoders. & Chapter~\ref{ch:courb} \\
    \addlinespace
    Check2HGI & Check-in-level representation: graph infomax extended to the check-in, so each visit carries its own vector. & Chapter~\ref{ch:mobiwac} \\
    \addlinespace
    Joint model & Cross-attention model on Check2HGI: one model for both tasks, next category and next region. & Chapter~\ref{ch:mobiwac} \\
    \bottomrule
  \end{tabularx}
\end{table}



\section{Multitask learning}
\label{sec:fund:mtl}

Multitask learning (MTL) trains related tasks together in the expectation that shared representations improve
generalization~\cite{caruana1997multitask}. This section explains how tasks share a model,
why that sharing may fail, and how multitask optimization responds to
the failure. This dissertation builds joint models for two evolving task pairs. The first pair combines static category
classification from a place embedding with next-category prediction from a check-in history. The final joint model
reads a check-in history to predict both the next category and the next region of a visit. Whether the shared
computation helps depends on the sharing topology and the training objective.

\subsection{Sharing topologies}

The first design choice is where task-specific computation ends and shared
computation begins.

\begin{definition}[Hard parameter sharing]\label{def:fund:hard}
Hard parameter sharing passes all inputs through a single shared trunk before
branching to task-specific output heads, so every task uses the same hidden
representations~\cite{ruder2017mtloverview}.
\end{definition}

\begin{definition}[Soft parameter sharing]\label{def:fund:soft}
Soft parameter sharing gives each task its own complete network and couples the
networks by penalizing differences between their
parameters~\cite{ruder2017mtloverview}.
\end{definition}

Definition~\ref{def:fund:hard} is the topology of the first joint model of this
dissertation, and it carries the risk this section addresses. A shared trunk need not
treat both tasks identically. Feature-wise Linear Modulation (FiLM) scales and shifts
features according to an additional input~\cite{perez2018film}, and MTLnet uses it to
condition its shared layers on task identity, so the two tasks read the same
parameters through different scalings.

\begin{definition}[Negative transfer]\label{def:fund:negtransfer}
Negative transfer occurs when joint training leaves a task worse than its dedicated
single-task model.
\end{definition}

Because hard sharing is rigid and soft sharing requires many parameters, several
architectures explore intermediate topologies. Cross-stitch units learn linear
combinations of activations from different task networks~\cite{misra2016cross}.
The multi-gate mixture-of-experts shares expert subnetworks but learns a
separate routing gate for each task~\cite{ma2018mmoe}. Progressive layered
extraction separates shared and task-specific experts across successive
gates~\cite{tang2020ple}, while DSelect-k makes expert selection sparse and
differentiable~\cite{hazimeh2021dselectk}. These methods selectively share
computation instead of forcing every task through one fixed trunk.

The joint model in Chapter~\ref{ch:mobiwac} builds a partial hard-shared interaction
subsystem based on this principle, using bidirectional cross-attention rather than
routing gates. The tasks do not reuse the same weight matrices, but their streams
cross within this jointly optimized module. Each stream queries the other to read its
features, which couples their gradients without forcing their representations into a
single bottleneck.

\subsection{Gradient conflict}

The second design choice depends on a quantity rather than on a preference: whether the
tasks are asking for the same update at all. Gradient conflict describes one source of
negative transfer, and it is a measured quantity rather than a figure of speech.

\begin{definition}[Gradient conflict]\label{def:fund:conflict}
Let $\mathbf{g}_i$ and $\mathbf{g}_j$ be the gradients of two task losses with respect
to the shared parameters, and let $\varphi_{ij}$ be the angle between them, so that
\begin{equation}\label{eq:fund:cosine}
    \cos\varphi_{ij}
    = \frac{\mathbf{g}_i^{\top}\mathbf{g}_j}
           {\lVert \mathbf{g}_i\rVert\,\lVert \mathbf{g}_j\rVert}.
\end{equation}
The two tasks conflict at that point when $\cos\varphi_{ij}<0$~\cite{yu2020pcgrad}.
\end{definition}

The cosine between their gradients reads how the two requested updates are aligned and
ignores how large they are: $+1$ means the tasks ask for the same update, $-1$ for
opposite ones, and $0$ that the requests are orthogonal, each leaving the other's
objective unchanged to first order. PCGrad changes a direction only when this cosine is
negative. Other methods may still react to differences in gradient magnitude or task
loss. In Chapter~\ref{ch:mobiwac}, two such methods exceed equal weighting on one
dataset but not on the others. Appendix~\ref{apx:cosine} evaluates the final joint
model and finds the task gradients equivalent to orthogonal within a margin of five
hundredths at every dataset measured there. A paired test separates the mean from zero
at three datasets, but those small departures remain within the equivalence margin.
The combined result is that the deviations are detectable but too small to represent
meaningful directional conflict at the scale considered by the balancers.

\subsection{Multi-objective optimization}

The third design choice is how several task losses determine one parameter update. A
measured conflict has to be acted on somewhere, and the place where it acts is the step
that turns two losses into one direction. Shared parameters can produce the negative
transfer of Definition~\ref{def:fund:negtransfer}, so task relationships must be measured
rather than assumed~\cite{standley2020tasks}. The training problem has
multiple losses, and a standard scalar objective writes them as
\begin{equation}\label{eq:fund:mtl-total}
    \mathcal{L}_{\mathrm{total}}(\boldsymbol{\theta})
    = \sum_{k=1}^{K} w_k\,\mathcal{L}_k(\boldsymbol{\theta}),
\end{equation}
over $K$ tasks and shared parameters $\boldsymbol{\theta}$, where $\mathcal{L}_k$ is the
loss of task $k$. A balancing method changes the weights $w_k$ or the update direction
that they induce. The scalar sum does not
remove the multi-objective nature of the problem~\cite{sener2018mgda}. One parameter
setting exhibits Pareto dominance over another when it is no worse on every task loss
and better on at least one. A setting is Pareto optimal when no other setting dominates
it, and the corresponding loss vectors form the Pareto front.

\subsubsection{Guarantees and their limits}

The guarantees of balancing methods are usually weaker than Pareto optimality. A
Pareto-stationary point is one at which a convex combination of the task gradients is
zero. This condition is necessary but not sufficient for Pareto
optimality~\cite{nash}. Nash-MTL proves convergence of a subsequence to such a point
and reaches Pareto optimality only under an additional convexity assumption that does
not hold for a deep network~\cite{nash}. CAGrad has Pareto-stationary fixed
points~\cite{liu2021cagrad}, and Aligned-MTL converges to one when task weights are
fixed in advance~\cite{senushkin2023aligned}. PCGrad makes no Pareto claim of its own. Under
conditions on two task gradients, loss curvature, and step size, it guarantees that
one projected update does not produce a higher multitask loss than the original
gradient~\cite{yu2020pcgrad}. Reaching the Pareto front would still leave the balance
between tasks unspecified~\cite{liu2021cagrad,senushkin2023aligned}. This
dissertation therefore claims no Pareto property for its models. It judges each task
against its dedicated single-task model under the tests defined in
Section~\ref{sec:fund:eval}.

\subsection{Loss-balancing methods}


Balancing methods differ in whether they modify loss weights or task-gradient
directions.

The first class sets the weights $w_k$ of Equation~\ref{eq:fund:mtl-total} from
different signals. Kendall et al.\ weight each loss by the task's homoscedastic
uncertainty~\cite{kendall2018uncertainty}. Chen et al.\ introduce GradNorm, which
balances gradient magnitudes and training rates~\cite{chen2018gradnorm}.

S. Liu et al.\ introduce Dynamic Weight Average alongside their attention
architecture, deriving each weight from the rate of change of that task's loss,
measured as the ratio of its two previous loss values~\cite{liu2019dwa}, and B. Liu et
al.\ propose FAMO, which tracks loss changes without computing every task gradient at
each step~\cite{liu2023famo}.

The second class changes the update direction instead. The multi-objective reading
begins with Sener and Koltun, who cast multitask learning as multi-objective
optimization~\cite{sener2018mgda}. Yu et al.\ introduce PCGrad, which projects away a
conflicting component in the sense of Definition~\ref{def:fund:conflict}~\cite{yu2020pcgrad}.
B. Liu et al.\ introduce CAGrad, which protects the worst per-task improvement within a
neighborhood of the average gradient~\cite{liu2021cagrad}. Navon et al.\ introduce
Nash-MTL, which treats gradient combination as a bargaining problem~\cite{nash}, and
Senushkin et al.\ introduce Aligned-MTL, which uses the condition number of the linear
system of task gradients as a stability criterion and aligns that system's orthogonal
components~\cite{senushkin2023aligned}.

These methods must be compared with
simple alternatives. Random loss weighting can be competitive~\cite{lin2022rlw};
specialized optimizers often fail to improve on a tuned fixed-weight
baseline~\cite{xin2022domtl}; and plain loss summation with standard regularization
can match or improve on them~\cite{kurin2022scalarization}. Surveys of general and
dense-prediction MTL summarize the same design
space~\cite{vandenhende2022mtl,yu2024survey}. For this dissertation, a balancing
method is useful only if it improves on a tuned fixed weighting.

\subsection{Multitask learning for mobility prediction}

Prior mobility applications show several ways to connect auxiliary targets to an
exact-place recommendation.

In mobility, MTL has been used almost entirely in the service of next place. MCARNN jointly
predicts activity and place through a context-aware recurrent
unit~\cite{Liao2018}. CSLSL instead predicts time, then activity, and then location in
a cascade~\cite{huang2024cslsl}. iMTL couples next-POI recommendation with the
handling of uncertain check-ins~\cite{Zhang2020}. Halder et al.\ develop multitask
next-POI recommenders that account for queue time and, in the journal extension, user
interest~\cite{Halder2021,Halder2022}. TME instead applies tree-guided multitask
embedding to static semantic POI annotation~\cite{Xu2023}. HAMTL sets location
prediction as its main task and category prediction as an auxiliary task, and its
hierarchical decoder predicts the category of the next destination first and then
refines the prediction of the specific location using that predicted
category~\cite{wang2025hamtl}. The same group had already published IeMTLF, an
interaction-enhanced multitask framework for next location
prediction~\cite{wang2024iemtlf}.

These methods confirm that category information and auxiliary tasks can support
mobility prediction, but they do not instantiate the pair studied in the final
chapter. Among the works reviewed here, none treats next category and next region as
co-equal end targets of one joint model.

This dissertation addresses that gap starting with MTLnet, which evaluates hard
parameter sharing over a static place embedding~\cite{silva2025mtlnet}. That initial
configuration fails to outperform dedicated single-task models, motivating the
representation and topology refinements developed in subsequent chapters.


\section{Datasets and evaluation}
\label{sec:fund:eval}

This section defines the datasets, metrics, reference points, and validation protocol
used to interpret the results. It also states which statistical tests support the
comparison verbs used in the final study.

\subsection{Datasets}

The geographic scope combines five state-level datasets with one city outside the
United States.

The experiments use two sources of check-ins. Gowalla is the main source, and the
dissertation evaluates Alabama, Arizona, Florida, California, and
Texas~\cite{cho2011gowalla}. Istanbul check-ins come from Massive-STEPS, which also
draws attention to the field's continued dependence on old
datasets~\cite{wongso2025massivesteps}. The two sources supply different administrative
units for the region target: a census tract in the five United States datasets, and a
\emph{mahalle}, a municipal neighborhood, in Istanbul. The Foursquare dataset for New
York and Tokyo
is another common benchmark~\cite{yang2015tsmc}, but it provides context only and is
not used in the experiments reported here.

\subsection{Metrics and reference points}

Each sequential target has its own metric, and the scores are not combined into one
claim of task performance.

\subsubsection{Next category}

Next-category prediction uses macro-F1 because the class distribution is imbalanced.
Food accounts for roughly one third of the check-ins in a representative state, so
plain accuracy can hide poor performance on less frequent classes. For $C$ classes,
macro-F1 is
\begin{equation}\label{eq:fund:macro-f1}
    \operatorname{MacroF1} = \frac{1}{C}\sum_{c=1}^{C}
    \frac{2P_cR_c}{P_c+R_c},
\end{equation}
where $P_c$ and $R_c$ are precision and recall for class $c$. The unweighted mean
gives every class equal importance~\cite{sokolova2009measures}. It does not show which
individual classes improve, and it can be low even when overall accuracy is high. The
models use unweighted cross-entropy rather than a reweighted loss. In
Chapter~\ref{ch:mobiwac}, class weighting lowered both category macro-F1 and region
accuracy.

\subsubsection{Next region}

Next-region prediction uses accuracy at 10 (Acc@10):
\begin{equation}\label{eq:fund:acc10}
    \operatorname{Acc@10} = \frac{1}{N}\sum_{i=1}^{N}
    \mathbb{I}\!\left[y_i \in \operatorname{Top10}(\hat{\mathbf{p}}_i)\right],
\end{equation}
where $y_i$ is the true region and $\operatorname{Top10}(\hat{\mathbf{p}}_i)$ contains
the ten regions with the largest predicted scores. The metric is the fraction of test
check-ins for which the correct region appears in this set. It does not distinguish
first place from tenth and does not measure the probability assigned to the true
region. Regions absent from the training partition count as errors. The reported
out-of-distribution discounted, or OOD-discounted, Acc@10 is therefore the Acc@10
measured on the in-distribution check-ins, multiplied by one minus the fraction of
check-ins whose region is out of distribution. Chapter~\ref{ch:mobiwac} gives the
complete evaluation procedure.

\subsubsection{Joint-model selection and floors}

Dedicated single-task models evaluate each target independently using its respective
metric. Evaluating a joint model, however, requires a single validation signal that
considers the performance of both outputs simultaneously. The validation protocol
achieves this by combining category macro-F1 and region Acc@10 into a scalar
selection score:
\begin{equation}\label{eq:fund:joint-selection}
    S_{\mathrm{joint}} = \sqrt{\operatorname{MacroF1} \times \operatorname{Acc@10}}.
\end{equation}
Under this \emph{joint-best} convention, both reported task scores come from the same
validation-selected checkpoint, reflecting the balanced performance of a single
deployable system.

To interpret these scores, model performance is evaluated between lower reference
floors and upper benchmark ceilings. At the lower bound, basic baselines establish
performance floors: a majority-class predictor for category classification, and a
first-order Markov model over training visits for sequential
transitions~\cite{gambs2012mmc}. At the upper bound, the primary reference ceiling for
evaluating multitask gains is the dedicated single-task model for each output.
Theoretical limits from prior work, such as the 93~percent potential predictability
reported for coarse location prediction~\cite{song2010limits}, do not serve as
ceilings for seven-class macro-F1 or region ranking.


\subsection{Preparation and data split}

How the check-ins are divided determines which population a result covers.

The studies use stratified cross-validation~\cite{kohavi1995crossval}, but the
split changes across chapters. Chapters~\ref{ch:cbic} and~\ref{ch:courb} stratify
samples, so one user's check-ins may occur in both training and validation.
Chapter~\ref{ch:mobiwac} instead uses a grouped, stratified splitter that keeps each
user on one side of a fold~\cite{pedregosa2011sklearn}. A seed in that chapter is one
complete repetition of the five-fold experiment with a different random
initialization. Four seeds over one fixed set of five folds produce twenty fitted
models per configuration. The fixed folds isolate random-initialization variation but
do not measure uncertainty from resampling the users.

\subsection{Comparison and statistical decisions}

The validation design determines which generalization claim and comparison verb each
result can support.

Chapters~\ref{ch:cbic} and~\ref{ch:courb} report fold means and standard deviations
without significance tests. The inferential rules therefore apply to
Chapter~\ref{ch:mobiwac}. Its superiority verdict uses a paired $t$ test on the four
per-seed means. A Wilcoxon signed-rank test on the twenty paired fold results is
reported alongside it and reaches the same decisions~\cite{wilcoxon1945}. Holm
correction controls family-wise error across the six dataset
comparisons~\cite{holm1979}. A superiority test cannot establish that a difference is
small enough to accept. For Alabama and Arizona, two one-sided tests establish that
the joint model is statistically non-inferior within a two-point
margin~\cite{lakens2017tost}. Accordingly, \emph{outperforms} is reserved for paired
superiority and \emph{matches} for statistical non-inferiority within that margin.


\section{Relevance}
\label{sec:fund:relevance}

The preceding sections introduce the prediction targets, the representations, multitask
learning, and the evaluation protocol one at a time. They are parts of a
single problem. This section states how they combine, before the three studies put
them to work.

\subsection{The parts of one problem}

The argument begins with the targets: Next place, next category, and next region have
different output spaces and different meanings for a correct answer. This dissertation
treats next category and next region as final outputs rather than as intermediate
steps toward next-place prediction.

Although this dissertation reports on those two targets, its three studies do not all
train the same pair. Chapters~\ref{ch:cbic}
and~\ref{ch:courb} pair next-category prediction with category classification, the
static task that labels a place from its own representation.
Chapter~\ref{ch:mobiwac} keeps next category but replaces that static task with next
region. The replacement follows from the representation rather than from a change of
interest: a place embedding gives a static classifier one fixed vector per place,
whereas a check-in-level representation describes a single visit, and the static task
is a less natural fit for that input than a second sequential target.

Two choices increase what the input representation must carry: treating next category
and next region as final outputs, and replacing the static task. A place embedding
can encode semantic and geographic relationships, but it assigns the same vector to
every visit to a place, so it carries nothing about the visit being predicted.
Check2HGI answers that limit by adding the visit as a level of its own.

Multitask learning is the mechanism that would let two tasks use what they have in
common, and it is not a free gain. Hard sharing can leave a task below its dedicated
single-task model. The balancing methods proposed to prevent that outcome often do not
improve on a tuned fixed-weight baseline. For the pair of tasks studied here, the
gradients are close to orthogonal on the shared parameters, which leaves such a method
little to correct. Whether joint training helps therefore cannot be assumed. It has to be
measured against the dedicated single-task model for each output, under a protocol
that does not favor the joint model.

Section~\ref{sec:fund:eval} defines that protocol, and the final study applies it in
full. Chapters~\ref{ch:cbic} and~\ref{ch:courb} instead use sample-stratified splits
and report fold means without significance tests. In the final study, user-disjoint
cross-validation limits leakage between a user's training and test visits. Macro-F1 and OOD-discounted
Acc@10 describe different outputs and are read against explicit reference points.
Paired superiority tests support the verb \emph{outperforms}, while non-inferiority
tests within a two-point margin support \emph{matches}. These rules prevent an
observed mean difference from being reported as a conclusion that its test does not
support.

\subsection{Gaps, questions, and answers}

Whether joint training helps this task pair is still open, and two gaps in the
literature explain why. The field has a place-level representation, but not a
check-in-level one built for these targets. It studies multitask learning, but it has
mostly treated the region as an intermediate signal on the way to a place rather than
as an end target.

Each study addresses one question. Chapter~\ref{ch:cbic} asks whether a
place-level representation and hard parameter sharing are sufficient for a joint
model. For that configuration, the joint model does not outperform the dedicated
single-task models. Chapter~\ref{ch:courb} keeps the model fixed, asks whether the
input representation, rather than the sharing topology, is the bottleneck, and finds
that it is. Chapter~\ref{ch:mobiwac} asks what a check-in-level representation and a
different sharing topology enable for two sequential targets. There, the joint model
outperforms the dedicated model on next category at all six datasets, and on next
region it outperforms at four of them and matches the dedicated model at the other
two. Chapter~\ref{ch:mobiwac} reports which datasets fall on each side.

Those answers also decide something practical. Anticipating what kind of place a
person will visit, and where that visit will occur, can support recommendation,
navigation, transit planning, and the allocation of resources by area. A service that
needs both predictions can now take them from a single model to maintain and one
forward pass, instead of two dedicated single-task models. The gain is operational,
not computational. Chapter~\ref{ch:mobiwac} reports the joint model as the larger
artifact, and a forward pass through it costs more than running the two dedicated
models, so the reduction is in the number of models to train and maintain.
